\chapter{Related Work}
\label{sec:chapterRelatedWork}

This chapter reviews the existing literature on machine learning-based fraud detection, with a focus on two publicly available benchmark datasets: the ULB Credit Card Fraud Detection dataset~\cite{ulb2013} and the Bank Account Fraud (BAF) suite~\cite{jesus2022turning}. The review is organised thematically rather than chronologically. We first survey the studies conducted on each dataset, highlighting the models, preprocessing choices, and evaluation protocols adopted. We then discuss cross-cutting concerns that emerge from the literature: the handling of extreme class imbalance, the role of deep learning in tabular fraud detection, the use of explainability techniques, and the choice of evaluation metrics and decision thresholds. The chapter closes with a summary of the identified research gaps that motivate the contributions of this thesis.

% ══════════════════════════════════════════════════════════════════════════
\section{Fraud Detection on the ULB Credit Card Dataset}
\label{sec:rw_ulb}
% ══════════════════════════════════════════════════════════════════════════

The ULB Credit Card Fraud Detection dataset, introduced by the Machine Learning Group of the Université Libre de Bruxelles, contains 284\,807 European credit card transactions collected over two days, of which only 492 ($\approx$0.17\%) are fraudulent~\cite{ulb2013}. All features except \texttt{Time} and \texttt{Amount} are anonymised principal components (V1--V28) obtained via PCA, which prevents domain-informed feature engineering but ensures privacy. Despite its age and relatively small size, the dataset remains a dominant benchmark in fraud detection research--—a reflection of the scarcity of public tabular datasets for this domain.

\paragraph{Georgieva et al.\ (2019)}
proposed a pattern-recognition neural network implemented in MATLAB, using a single hidden layer of 13 neurons trained with scaled conjugate gradient backpropagation~\cite{georgieva2019neural}. To address class imbalance, the authors applied random undersampling to a balanced subset of 900 instances (450 per class). On this reduced set, the model achieved an $F_1$ score of 94.40\% and a Matthews Correlation Coefficient (MCC) of 88.62\%. However, the heavy undersampling discards over 99.6\% of the legitimate transactions, yielding a training distribution that is unrepresentative of the operational class ratio. Furthermore, the study reports metrics computed on the balanced subset rather than on the full imbalanced test set, which limits the practical significance of the results.

\paragraph{AlEmad (2022)}
compared K-Nearest Neighbours ($K$\,=\,3 and 7), Naïve Bayes, Logistic Regression, and Support Vector Machine on the ULB dataset using a 70:30 holdout split with no balancing strategy~\cite{alemad2022credit}. Performance was evaluated exclusively via accuracy, with the best-performing model (SVM) achieving 99.94\%. While this value appears high, it is misleading: a trivial classifier that predicts all transactions as legitimate would achieve $\approx$99.83\% accuracy on this dataset. The absence of precision, recall, $F_1$, PR-AUC, or any threshold-aware metric means the study provides virtually no insight into the models' ability to detect the rare fraudulent class.

\paragraph{Mosa et al.\ (2024)}
combined 15 meta-heuristic optimisation (MHO) algorithms with 9 transfer functions for binary feature selection, applied to Random Forest and SVM classifiers~\cite{mosa2024ccfd}. The study employed undersampling to a balanced set of 984 instances (492 per class) and reported accuracy as the primary metric. The best combination (SFO-RF with V-shaped transfer function V1) achieved 97.79\% accuracy on the balanced subset, improving on the non-optimised RF baseline by approximately 4 percentage points. The work demonstrates the viability of MHO-based feature selection on high-dimensional PCA features; however, the severe undersampling and reliance on accuracy as the sole metric are notable limitations. As a robustness check, the authors evaluated on the full imbalanced 20\% test set (284\,807 instances) and reported 97.14\% accuracy, though no class-specific metrics were provided for this evaluation.

\paragraph{Du et al.\ (2024)}
proposed a multi-stage pipeline combining an autoencoder for representation learning, XGBoost for probabilistic scoring, and SMOTE augmented by a Conditional Generative Adversarial Network (CGAN) for synthetic minority generation~\cite{du2024novel}. A distinguishing feature of this work is its explicit treatment of the decision threshold: the authors swept $\tau$ from 0 to 1 in steps of 0.01 and selected the threshold that maximised MCC. At $\tau = 0.35$, the pipeline achieved MCC\,=\,0.8845 and TPR\,=\,0.7839; at $\tau = 0.05$, TPR increased to 0.8929 at the cost of lower MCC (0.773). This threshold-aware evaluation is rare in the ULB literature and represents a methodological strength. The study used a stratified 60:20:20 split, and SMOTE was applied only to the training partition, avoiding leakage.

\paragraph{Wu et al.\ (2025)}
introduced a Continuous-Coupled Neural Network (CCNN) inspired by pulse-coupled neural networks from computational neuroscience~\cite{wu2025ccnn}. SMOTE was applied to upsample the minority class from 492 to 284\,315 instances. At the best temporal parameter ($T$\,=\,2), the authors reported accuracy of 0.9998, precision of 0.9996, recall of 1.0000, and $F_1$ of 0.9998. These near-perfect metrics warrant scrutiny. Critically, the description of the experimental protocol indicates that SMOTE was applied to the entire dataset \emph{before} the 70:30 train/test split, meaning that synthetic samples generated from training-set fraudulent transactions may appear in the test set. This constitutes data leakage, as the model is effectively evaluated on data derived from its own training distribution. When leakage is present, inflated metrics are expected and the results cannot be meaningfully compared with those of correctly designed experiments. The study also reports only accuracy-family metrics and does not evaluate PR-AUC or MCC.

\paragraph{Jansson and Axelsson (2020)}
investigated the application of Federated Averaging~\cite{mcmahan2017communication} to credit card fraud detection, using a centrally trained MLP as the baseline and comparing it against federated configurations with varying client fractions ($C \in \{0.5, 0.75, 1.0\}$) and local epochs ($E \in \{2, 3\}$)~\cite{jansson2020federated}. The authors adopted the area under the Precision--Recall curve (AUPRC) as the primary metric, arguing that ROC-AUC is inappropriate for severely imbalanced settings---a position supported by Saito and Rehmsmeier~\cite{saito2015precision}. The best federated model achieved AUPRC\,=\,0.868 versus 0.823 for the centralised baseline, suggesting that label heterogeneity across clients can act as an implicit regulariser.

An important finding concerns the effect of SMOTE in a federated setting. The authors experimented with oversampling at different ratios (1\%, 10\%, 50\%) applied exclusively on each client's local training partition. Aggressive oversampling (10\% and 50\%) proved catastrophic, generating thousands of false positives because the synthetic fraudulent samples were too similar to legitimate transactions. Even at 1\% upsampling, the model exhibited signs of overfitting to the minority class: while the number of true positives increased, AUPRC decreased, indicating that precision degraded faster than recall improved. The authors concluded that \emph{``the synthetic fraudulent datapoints are too similar to the non-fraudulent transactions''}~\cite{jansson2020federated}. This finding is particularly relevant to the present work, as it corroborates the observation---common in the broader literature on imbalanced learning---that oversampling can be counterproductive when the class overlap is high.

\paragraph{Summary.}
Table~\ref{tab:rw_ulb_summary} summarises the methodological choices across the six ULB studies. The landscape is characterised by considerable heterogeneity in evaluation protocols: metrics range from accuracy alone to AUPRC, balancing strategies vary from no treatment to aggressive SMOTE, and only one study~\cite{du2024novel} explicitly tunes the decision threshold.

\begin{table}[ht]
\centering
\caption{Summary of related work on the ULB Credit Card 2013 dataset.}
\label{tab:rw_ulb_summary}
\small
\begin{tabular}{l l l l c}
\toprule
Study & Best Model & Balancing & Primary Metric & Threshold \\
\midrule
Georgieva et al.\ (2019) & ANN          & RUS (900)       & $F_1$ = 0.944         & No  \\
AlEmad (2022)             & SVM          & None            & Acc = 0.999           & No  \\
Mosa et al.\ (2024)      & SFO-RF       & RUS (984)       & Acc = 0.978           & No  \\
Du et al.\ (2024)        & AE-XGB-CGAN  & SMOTE + CGAN    & MCC = 0.885           & Yes \\
Wu et al.\ (2025)        & CCNN         & SMOTE (leaked)  & $F_1$ = 0.9998$^{*}$  & No  \\
Jansson \& Axelsson (2020)& FedAvg MLP  & SMOTE (1\%)     & AUPRC = 0.868         & No  \\
\bottomrule
\multicolumn{5}{l}{\footnotesize $^{*}$Likely inflated due to SMOTE applied before train/test split (data leakage).}
\end{tabular}
\end{table}


% ══════════════════════════════════════════════════════════════════════════
\section{Fraud Detection on the BAF Dataset}
\label{sec:rw_baf}
% ══════════════════════════════════════════════════════════════════════════

The Bank Account Fraud (BAF) dataset suite was introduced at NeurIPS 2022 by Jesus et al.~\cite{jesus2022turning}. It consists of a Base dataset and five variants (I--V), each containing approximately one million synthetic-but-realistic instances generated via CTGAN from a real-world bank account opening fraud detection system. The fraud prevalence is approximately 1.1\%, and the feature space comprises 30 interpretable attributes (including demographic, behavioural, and device-related variables). Unlike the ULB dataset, the BAF features are not anonymised, enabling domain-informed analysis and explainability studies.

\paragraph{Jesus et al.\ (2022)}
introduced the BAF suite as a benchmark for evaluating ML methods under realistic conditions, with an emphasis on fairness~\cite{jesus2022turning}. The six dataset variants differ systematically in group size disparity, prevalence disparity, and separability disparity across a protected attribute (customer age $\leq$\,50 vs.\ $>$\,50), allowing controlled study of bias--performance trade-offs. A temporal train/test split (first 6 months for training, last 2 for testing) mirrors realistic deployment, where models must generalise to future data.

The authors trained 100 configurations of LightGBM on each variant and measured True Positive Rate at 5\% False Positive Rate (TPR@5\%FPR), a metric they describe as \emph{``typically imposed by clients in the fraud detection domain.''} Predictive equality (ratio of FPRs between demographic groups) was used as the fairness metric. A key finding is that the top-performing models on the Base dataset are not necessarily the best on the temporal variants (IV and V), confirming that \emph{``performant models in static environments may fall short in more realistic, dynamic ones''}~\cite{jesus2022turning}. The dataset generation is well-documented, though the authors acknowledge that differential privacy guarantees are not formally proven and that the data is subject to selective labelling (only accepted applications receive ground-truth labels).

\paragraph{Sun et al.\ (2025)}
used the BAF Base variant to investigate the thesis that, under extreme imbalance, \emph{``objective choice and operating point matter at least as much as architecture''}~\cite{sun2025objective}. The authors compared Logistic Regression, SVM (RBF kernel), Random Forest, LightGBM, and a GRU network. Classical models were tuned via grid search on a 10\,000-instance development set, using $F_1$ weighted as the selection criterion and including \texttt{class\_weight='balanced'} in the search space. The GRU used inverse-frequency class weights in its cross-entropy loss.

The results reveal a stark failure mode: all four classical models achieve $\geq$0.98 weighted $F_1$ but only 1--6\% fraud recall, meaning they effectively learn to classify all instances as non-fraudulent. Even with \texttt{class\_weight='balanced'} enabled, the $f_1$-weighted selection criterion drives convergence towards majority-class solutions. The GRU, by contrast, achieves 78\% fraud recall but at only 5\% precision, generating approximately 19 false positives per true positive. The authors' feature importance analysis shows that all five models agree on the top predictive features (\texttt{name\_email\_similarity}, \texttt{velocity\_24h}, \texttt{velocity\_6h}), demonstrating that ``the failure of the classical models was not one of feature identification, but one of objective''~\cite{sun2025objective}. The study uses a fixed threshold of $\tau$\,=\,0.5 and acknowledges this is suboptimal, recommending analysis of the full PR curve as a next step.

\paragraph{AbouGrad and Sankuru (2025)}
proposed a decentralised anomaly detection framework using deep autoencoders, applied to both the ULB 2013 dataset and the full BAF suite~\cite{abougrad2025online}. Each dataset was treated as an independent institutional node with no parameter sharing, in a design motivated by GDPR compliance. The autoencoder (encoder: 16$\rightarrow$8$\rightarrow$4 neurons; decoder: symmetric) was trained exclusively on legitimate transactions, with anomalies flagged when the reconstruction error exceeded the 98th percentile of the training distribution. The achieved ROC-AUC of 0.662 is modest, and the PR-AUC of 0.012 reflects the inherent difficulty of unsupervised anomaly detection under extreme class imbalance. The authors conducted a threshold sensitivity analysis comparing the 95th, 98th, and 99th percentiles, and observed the expected precision--recall trade-off. While the study is methodologically sound in avoiding leakage, the strictly unsupervised approach---which ignores available labels---leaves substantial performance on the table compared to supervised alternatives.

\paragraph{Alves et al.\ (2025)}
extended the BAF Base dataset to create FiFAR---a benchmark for Learning to Defer (L2D) in human-AI decision-making~\cite{alves2025benchmarking}. The framework generates 50 synthetic fraud analysts with instance-dependent noise, varying consistency (measured by Cohen's $\kappa$), and bias against older customers. A LightGBM alert model is trained on the first three months and validated on the fourth, optimising recall at 5\% FPR via the Neyman--Pearson criterion, which yields an explicit cost ratio $\lambda = (1-t)/t \approx 0.057$ between false positive and false negative costs. The L2D experiments compare two deferral algorithms---One-vs-All (OvA) and DeCCaF---against baselines (full rejection, classifier-only, random assignment). Both L2D methods reduce total misclassification cost by 11--15\% relative to random assignment, but, importantly, \emph{``performance rankings of L2D algorithms vary significantly based on the available experts''}~\cite{alves2025benchmarking}. While the L2D paradigm is outside the scope of this thesis, the study's rigorous use of temporal splits, cost-sensitive thresholding, and the Neyman--Pearson formulation provides a reference for principled fraud detection evaluation.

\paragraph{Summary.}
Table~\ref{tab:rw_baf_summary} summarises the BAF-related studies. The BAF literature benefits from having been established more recently (2022), which is reflected in generally higher methodological rigour: temporal splits, threshold-aware metrics, and explicit fairness considerations are more common than in the ULB literature.

\begin{table}[ht]
\centering
\caption{Summary of related work on the BAF dataset suite.}
\label{tab:rw_baf_summary}
\small
\begin{tabular}{l l l l c}
\toprule
Study & Best Model & Balancing & Primary Metric & Threshold \\
\midrule
Jesus et al.\ (2022)  & LightGBM    & None          & TPR@5\%FPR        & 5\% FPR  \\
Sun et al.\ (2025)    & GRU         & Class weights & Fraud Recall      & Fixed 0.5 \\
AbouGrad \& S.\ (2025) & Autoencoder & Unsupervised  & ROC-AUC = 0.662   & 98th pct. \\
Alves et al.\ (2025)  & LightGBM    & Cost-weighted & Cost / Recall@5\% & Neyman--Pearson \\
\bottomrule
\end{tabular}
\end{table}


% ══════════════════════════════════════════════════════════════════════════
\section{Class Imbalance Handling in Fraud Detection}
\label{sec:rw_imbalance}
% ══════════════════════════════════════════════════════════════════════════

Extreme class imbalance is the defining characteristic of fraud detection datasets, with fraudulent instances typically comprising less than 2\% of the data. The literature reveals three main families of approaches: data-level resampling, algorithm-level cost-sensitive learning, and hybrid methods.

\paragraph{Random undersampling.}
Several ULB studies employ random undersampling (RUS) to create balanced training sets~\cite{georgieva2019neural, mosa2024ccfd}. While this equalises the class distribution, severe undersampling---such as reducing the dataset from 284\,807 to 900 or 984 instances---discards the vast majority of the legitimate class and yields a training set that is unrepresentative of the deployment distribution. Metrics computed on the balanced subset can appear high (e.g., $F_1 > 0.94$), but these are not indicative of performance in the original imbalanced setting.

\paragraph{SMOTE and its variants.}
SMOTE~\cite{chawla2002smote} generates synthetic minority samples by interpolating between existing minority instances and their nearest neighbours. Wu et al.~\cite{wu2025ccnn} used SMOTE to upsample the ULB minority class from 492 to 284\,315 instances, but applied it before the train/test split, introducing leakage. Du et al.~\cite{du2024novel} combined SMOTE with a CGAN to produce more diverse synthetic fraudulent samples and applied both techniques correctly on the training partition only.

The most informative study on the effects of SMOTE in fraud detection is that of Jansson and Axelsson~\cite{jansson2020federated}, who evaluated multiple oversampling ratios (1\%, 10\%, 50\%) in a federated setting. Their results demonstrate a consistent pattern: aggressive oversampling degraded performance by introducing synthetic fraud samples that overlapped with the legitimate class distribution, resulting in a rise of false positives. Even the most conservative ratio (1\%) caused the model to overfit to the upsampled minority class, worsening AUPRC despite improving the raw true positive count. These findings challenge the common assumption that oversampling unconditionally benefits minority-class detection.

\paragraph{Cost-sensitive learning.}
Sun et al.~\cite{sun2025objective} and Alves et al.~\cite{alves2025benchmarking} use class weighting as an alternative to resampling. Sun et al.\ included \texttt{class\_weight='balanced'} in the hyperparameter grid for classical models and used inverse-frequency weights for the GRU. Notably, enabling class weighting within standard libraries proved insufficient to overcome the majority-class bias when the global selection metric ($f_1$-weighted) did not explicitly prioritise minority-class recall. Alves et al.\ took a cost-sensitive approach grounded in the Neyman--Pearson framework, deriving an asymmetric cost ratio $\lambda$ from the operating threshold.

\paragraph{Implication.}
The literature suggests that no single imbalance-handling strategy is universally beneficial. The effectiveness of a strategy depends on the interaction between the resampling technique, the model family, and the evaluation metric. This observation motivates the full factorial analysis conducted in this thesis (Section~\ref{sec:imbalance_strategies}), where seven strategies are systematically crossed with five models and multiple datasets.


% ══════════════════════════════════════════════════════════════════════════
\section{Deep Learning for Tabular Fraud Detection}
\label{sec:rw_deep_learning}
% ══════════════════════════════════════════════════════════════════════════

Despite the success of deep learning in computer vision and natural language processing, its application to tabular fraud detection remains limited and inconclusive.

Among the studies reviewed, deep learning architectures appear in several forms: a single-hidden-layer ANN~\cite{georgieva2019neural}, a brain-inspired CCNN~\cite{wu2025ccnn}, an autoencoder for anomaly detection~\cite{du2024novel, abougrad2025online}, and a GRU network~\cite{sun2025objective}. With the exception of the GRU (which was specifically designed to test objective-choice hypotheses), none of these architectures represents a modern, purpose-built tabular deep learning model. Shallow feedforward networks~\cite{georgieva2019neural, jansson2020federated} and domain-specific architectures~\cite{wu2025ccnn} are not comparable to recent advances such as FT-Transformer~\cite{gorishniy2021revisiting}, TabNet~\cite{arik2021tabnet}, or TabTransformer~\cite{huang2020tabtransformer}, which leverage self-attention mechanisms to model feature interactions in tabular data.

Notably, no study in the reviewed literature applies a Transformer-based model to either the ULB or BAF datasets. Sun et al.~\cite{sun2025objective} explicitly suggest benchmarking against TabTransformer as future work. This gap is significant given that recent work on tabular benchmarks has shown that FT-Transformer and similar architectures can match or surpass gradient-boosted trees on certain datasets~\cite{gorishniy2021revisiting}, particularly when feature interactions are complex. Whether this advantage extends to the specific characteristics of fraud detection---extreme imbalance, high class overlap, and temporal dynamics---remains an open question that this thesis addresses by including an FT-Transformer in the model comparison (Section~\ref{sec:models}).


% ══════════════════════════════════════════════════════════════════════════
\section{Explainability in Fraud Detection}
\label{sec:rw_xai}
% ══════════════════════════════════════════════════════════════════════════

Explainable AI (XAI) is frequently mentioned in the fraud detection literature as a desirable goal, yet its substantive application remains rare in the studies reviewed.

AbouGrad and Sankuru~\cite{abougrad2025online} dedicate a section to XAI in financial fraud detection, discussing SHAP~\cite{lundberg2017unified} and LIME~\cite{ribeiro2016why} at a conceptual level and recommending their integration as future work. Sun et al.~\cite{sun2025objective} compute feature importance across all five model families (via coefficients, Gini importance, gain, and permutation importance) and observe that all models agree on the top predictive features, providing useful insight into the signal structure of the BAF dataset. Alves et al.~\cite{alves2025benchmarking} analyse feature dependence in both their synthetic expert models and the LightGBM alert model, identifying that the model score is the most important feature driving expert decisions.

However, none of the ten studies applies post-hoc explanation methods (such as SHAP or LIME) at the instance level to analyse individual predictions, nor does any study investigate the trade-off between model performance and interpretability---for example, whether a more interpretable model (such as Logistic Regression) sacrifices detection power relative to a black-box alternative (such as LightGBM or a neural network), and whether the lost performance can be recovered through threshold tuning or resampling.

This gap is particularly notable for the BAF dataset, where the non-anonymised feature space enables semantically meaningful explanations. The present thesis addresses this by including SHAP-based global and local explanations for all model families on the BAF suite (Section~\ref{sec:interpretability}) and by analysing whether the most important features are consistent across models and across dataset variants.


% ══════════════════════════════════════════════════════════════════════════
\section{Evaluation Protocols and Metrics}
\label{sec:rw_evaluation}
% ══════════════════════════════════════════════════════════════════════════

A recurring concern in the fraud detection literature is the choice of evaluation metrics and the rigour of experimental protocols. Three issues deserve attention.

\paragraph{Accuracy as a metric under extreme imbalance.}
Several studies report accuracy as their primary or sole metric~\cite{alemad2022credit, mosa2024ccfd}. Under the class distributions typical of fraud detection (0.17\% fraud in ULB, 1.1\% in BAF), a trivial all-negative classifier achieves $>$99\% accuracy, rendering the metric uninformative for assessing fraud detection capability. The studies that adopt ranking-based metrics---AUPRC~\cite{jansson2020federated}, MCC~\cite{du2024novel}, or TPR at a fixed FPR~\cite{jesus2022turning, alves2025benchmarking}---provide substantially more insight into model performance on the minority class.

\paragraph{Threshold selection and leakage.}
The decision threshold is a critical but often neglected component of the evaluation pipeline. Among the ten studies reviewed, only Du et al.~\cite{du2024novel} perform an explicit threshold sweep, selecting $\tau$ based on MCC maximisation. Jesus et al.~\cite{jesus2022turning} and Alves et al.~\cite{alves2025benchmarking} bypass the threshold choice by reporting TPR at a fixed FPR, which provides an operating-point-independent comparison. Sun et al.~\cite{sun2025objective} use the default $\tau = 0.5$ and acknowledge this is suboptimal. The remaining studies either do not discuss threshold selection or implicitly use the default, which is known to be inappropriate for heavily skewed class distributions~\cite{he2009learning}.

Beyond threshold choice, the application of preprocessing steps on the entire dataset before the train/test split---as observed in Wu et al.~\cite{wu2025ccnn} for SMOTE---constitutes data leakage that invalidates the reported metrics. In the broader fraud detection literature, this is a documented concern: models with leakage may report near-perfect metrics that do not generalise to unseen data.

\paragraph{Temporal evaluation.}
Fraud detection is inherently temporal: models trained on historical data must generalise to future transactions, and fraudster behaviour evolves over time. Among the reviewed studies, only the BAF-based works~\cite{jesus2022turning, alves2025benchmarking} employ temporal train/test splits (training on earlier months, testing on later months). The ULB studies uniformly use random or stratified splits, which may overestimate performance by allowing temporal information to leak between training and test sets.


% ══════════════════════════════════════════════════════════════════════════
\section{Summary and Research Gaps}
\label{sec:rw_gaps}
% ══════════════════════════════════════════════════════════════════════════

The reviewed literature reveals several gaps that the present thesis aims to address:

\begin{enumerate}
    \item \textbf{Inconsistent evaluation protocols.} The ULB studies vary widely in their choice of metrics, balancing strategies, and splitting procedures, making cross-study comparison unreliable. This thesis adopts a unified, leakage-free evaluation protocol (Section~\ref{sec:protocol}) applied consistently across all models, strategies, and datasets.

    \item \textbf{Lack of systematic imbalance-handling analysis.} While individual studies test one or two balancing strategies, no work conducts a full factorial comparison of multiple strategies across multiple models and datasets. The negative effects of SMOTE documented by Jansson and Axelsson~\cite{jansson2020federated} and the limitations of class weighting observed by Sun et al.~\cite{sun2025objective} suggest that the interaction between strategy and model is non-trivial. This thesis evaluates seven strategies $\times$ five supervised models $\times$ multiple datasets (Section~\ref{sec:imbalance_strategies}).

    \item \textbf{Absence of modern tabular deep learning.} No reviewed study applies Transformer-based tabular architectures (FT-Transformer, TabNet, TabTransformer) to either dataset. This thesis includes FT-Transformer as a model family, testing whether attention-based architectures offer advantages over gradient-boosted trees in the specific context of fraud detection under extreme imbalance.

    \item \textbf{No substantive XAI application.} While explainability is discussed conceptually, no study provides instance-level explanations or analyses the performance--interpretability trade-off. This thesis applies SHAP-based explanations to all model families and examines feature importance consistency across models and dataset variants (Section~\ref{sec:interpretability}).

    \item \textbf{Threshold selection as an afterthought.} Only one ULB study and two BAF studies address threshold selection rigorously. This thesis treats threshold selection as a first-class experimental variable, comparing four threshold strategies and measuring their impact on model rankings and operational metrics (Section~\ref{sec:threshold_study}).
\end{enumerate}

Together, these contributions aim to provide a comprehensive, methodologically rigorous benchmark for fraud detection under extreme class imbalance, bridging the gaps identified in the current literature.
